% =====================================================================
%  5.2 Arealsetningen — fasit
% =====================================================================

\fk{Del 1 — Arealsetningen og beviset}
\fo{1.1} a) $6$ \ b) $14$ \ c) $5$ \ d) $6\sqrt2$
\fo{1.2} a) Vinkel $B$ \ b) $T = \frac12 \cdot 6 \cdot 8 \cdot \sin 30^\circ = 12$
\fo{1.3} a) $31{,}0$ cm$^2$ \ b) $47{,}1$ cm$^2$ (vinkelen $B$ ligger mellom $a$ og $c$)
\fo{1.4} a) $27\sqrt3$ \ b) $10\sqrt2$ \ c) $15\sqrt3$
\fo{1.5} Feilen er overgangen til $24\cdot(-\frac{\sqrt3}{2})$: $\sin 120^\circ = \sin 60^\circ = \frac{\sqrt3}{2}$, ikke negativ. Riktig: $T = 24 \cdot \frac{\sqrt3}{2} = 12\sqrt3$. Et areal kan ikke være negativt.
\fo{1.6} a) $\frac12 \cdot 4{,}2 \cdot 3{,}5 \cdot \sin 75^\circ \approx 7{,}1$ m$^2$ \ b) $\approx 7{,}1$ m$^2$. Samme areal, fordi $\sin 105^\circ = \sin 75^\circ$.
\fo{1.7} a) $\angle CBD = 180^\circ - B$ \ b) $h = a\sin(\angle CBD)$ \ c) $h = a\sin(180^\circ - B) = a\sin B$ \ d) Grunnlinje $c$ og høyde $h$ gir $T = \frac12 c h = \frac12 ac\sin B$.
\fo{1.8} a) Noen ganger. Med $C = 30^\circ$ og $C = 150^\circ$ er arealet $20 \cdot \frac12 = 10$ begge ganger, men med $C = 60^\circ$ er det $10\sqrt3$.
b) Noen ganger. Arealet øker fram til $C = 90^\circ$, der $\sin C = 1$ er størst, og minker deretter.
\fo{1.9} Eksempel: $a = b = 4$. Med $C = 45^\circ$ er $T = 8 \cdot \frac{\sqrt2}{2} = 4\sqrt2$. Med $C = 90^\circ$ er $T = 8$. Men $2 \cdot 4\sqrt2 = 8\sqrt2 \neq 8$. Andre moteksempler er riktige.
\fo{1.10} a) Alle vinkler er $60^\circ$ og alle sider $s$: $T = \frac12 s \cdot s \cdot \sin 60^\circ = \frac{\sqrt3}{4}s^2$. \ b) $\frac{\sqrt3}{4}s^2 = 16\sqrt3$ gir $s^2 = 64$, så $s = 8$.
\fo{Sjekk 1} a) $\frac12 \cdot 9 \cdot 4 \cdot \frac12 = 9$ \ b) Faktoren $\frac12$ mangler. Riktig er $T = \frac12 ab\sin C$. \ c) $\sin 90^\circ = 1$, så $T = \frac12 ab \cdot 1$. Katetene er grunnlinje og høyde.

\fk{Del 2 — Bruke arealsetningen}
\fo{2.1} a) $a = 5$ \ b) $a = 8$ \ c) $a = 6$
\fo{2.2} a) $20 \cdot \sin C = 10$ og $\sin C = \frac12$ \ b) $C = 30^\circ$ eller $C = 150^\circ$
\fo{2.3} a) $36{,}9^\circ$ og $143{,}1^\circ$ \ b) $14{,}5^\circ$ og $165{,}5^\circ$ \ c) $64{,}2^\circ$ og $115{,}8^\circ$
\fo{2.4} a) $\sin C = \frac{2 \cdot 28}{9 \cdot 7} \approx 0{,}889$, så $C \approx 62{,}7^\circ$ eller $C \approx 117{,}3^\circ$.
b) $\frac12 \cdot 9 \cdot 7 \cdot \sin 117{,}3^\circ \approx 28{,}0$
\fo{2.5} a) $\sin C = \frac{24\sqrt2}{48} = \frac{\sqrt2}{2}$, så $C = 45^\circ$ eller $C = 135^\circ$.
b) $\sin C = \frac{28}{28} = 1$, så bare $C = 90^\circ$.
c) $\sin C = \frac{30}{24} = \frac54 > 1$. Ingen trekant. Det største arealet med $a = 6$ og $b = 4$ er $12$.
\fo{2.6} $T_{ABC} = \frac12 \cdot 9 \cdot 12 \cdot \sin 35^\circ \approx 31{,}0$ og $T_{ACD} = \frac12 \cdot 12 \cdot 8 \cdot \sin 50^\circ \approx 36{,}8$. Sammen $\approx 67{,}7$ m$^2$.
\fo{2.7} Feilen er siste steg: $\sin C = \frac{\sqrt3}{2}$ gir både $C = 60^\circ$ og $C = 120^\circ$. Begge gir areal $12\sqrt3$, så det finnes to trekanter.
\fo{2.8} a) Eksempel: $a = 6$, $b = 8$: $\frac12 \cdot 6 \cdot 8 \cdot \frac12 = 12$. Flere svar er riktige.
b) Uendelig mange. Alle $a$ og $b$ med $\frac12 ab \cdot \frac12 = 12$, altså $ab = 48$, gir en trekant (for eksempel $4$ og $12$).
c) $\sin C = \frac{24}{24} = 1$, så $C = 90^\circ$. Nøyaktig én trekant.
\fo{2.9} $\sin C = \frac{2T}{60}$. To trekanter: $0 < T < 30$ (da er $\sin C < 1$). Én trekant: $T = 30$ ($\sin C = 1$, $C = 90^\circ$). Ingen: $T > 30$, fordi $\sin C$ ikke kan være større enn $1$. Største areal er $\frac12 \cdot 6 \cdot 10 = 30$.
\fo{Sjekk 2} a) $2{,}25a = 18$, så $a = 8$ \ b) $17{,}5^\circ$ og $162{,}5^\circ$ \ c) Fordi $\sin(180^\circ - v) = \sin v$. Da er $\frac12 ab\sin v$ det samme for begge.

\fk{Ekstra utfordring}
\fo{U1} Hver trekant har to sider på $4$ og vinkel $60^\circ$: $\frac12 \cdot 4 \cdot 4 \cdot \frac{\sqrt3}{2} = 4\sqrt3$. Sekskanten: $6 \cdot 4\sqrt3 = 24\sqrt3$.
\fo{U2} a) $\sin C = \frac{3200}{4000} = 0{,}8$, så $C \approx 53{,}1^\circ$ eller $C \approx 126{,}9^\circ$. b) To beiter. Begge har sidene $50$ m og $80$ m og samme høyde, og dermed samme areal.
\fo{U3} Arealet er $\frac12 \cdot 7 \cdot 12 \cdot \sin C = 42\sin C$. Siden $\sin C \le 1$, er det størst når $\sin C = 1$, altså $C = 90^\circ$. Største areal er $42$.
\fo{U4} Diagonalen deler parallellogrammet i to like trekanter med sidene $a$ og $b$ og vinkelen $v$ mellom (motstående vinkler er like, og sidene er parvis like). Hver har areal $\frac12 ab\sin v$. Sammen $ab\sin v$. For stump $v$ gjelder arealsetningen også, så svaret er det samme.
\fo{U5} a) Trekantene har to sider på $1$ (radiene) og vinkelen $\frac{360^\circ}{n}$ mellom. Hver har areal $\frac12 \cdot 1 \cdot 1 \cdot \sin\frac{360^\circ}{n}$. Med $n$ trekanter: $T_n = \frac n2\sin\frac{360^\circ}{n}$.
b) $T_4 = 2\sin 90^\circ = 2$, $T_6 = 3\sin 60^\circ = \frac{3\sqrt3}{2}$, $T_{12} = 6\sin 30^\circ = 3$.
\fo{U6} a) For $n = 6$: $\frac{3\sqrt3}{2} \approx 2{,}598$. For $n = 12$: $3$ (programmet skriver omtrent $3{,}0$). b) $T_n$ nærmer seg $\pi \approx 3{,}1416$ ($T_{1536} \approx 3{,}14158$).
c) Når $n$ vokser, fyller $n$-kanten sirkelen bedre og bedre. Sirkelen med radius $1$ har areal $\pi \cdot 1^2 = \pi$.
d) Funksjonen \lstinline[style=laereverk,numbers=none,frame=none]|math.sin| regner i radianer. \lstinline[style=laereverk,numbers=none,frame=none]|math.radians| gjør grader om til radianer.
